% ch2.tex -- Module 5, Chapter 2: Circular Motion: Forces and Real Situations
% Vertical circles (forces at the top and bottom, energy) are taught in full in ch3.
\chapter[Circular Motion: Forces]{Circular Motion: Forces and Real Situations}\label{ch2:chap}
\covers{M5.IQ2.01, M5.IQ2.02, M5.IQ2.03}

\inquiry{Why do objects move in circles?}

\begin{goals}
  \item explain why an object moving in a circle at a constant speed is accelerating
  \item solve problems with $v$, $a_\mathrm{c}$ and $F_\mathrm{c}$, and name the real force
        that provides $F_\mathrm{c}$
  \item plan and evaluate an investigation of $F_\mathrm{c}$, $m$, $v$ and $r$
\end{goals}

\begin{syllabus}
  \item M5.IQ2.01: investigate how centripetal force, mass, speed and radius are related
  \item M5.IQ2.02: analyse the forces on objects in uniform circular motion
  \item M5.IQ2.03: solve problems with $a_\mathrm{c} = v^2/r$, $v = 2\pi r/T$,
        $F_\mathrm{c} = mv^2/r$ and $\omega = \Delta\theta/t$
\end{syllabus}

%% ==================================================================
\section{Uniform circular motion and centripetal force}\label{ch2:sec:ucm}
\covers{M5.IQ2.01, M5.IQ2.02, M5.IQ2.03}

\begin{bigidea}
An object moving in a circle at a constant speed is accelerating, because its direction
keeps changing. The acceleration points to the centre. A net force towards the centre,
the centripetal force, causes it.
\end{bigidea}

\begin{keywords}
  \kw{Uniform circular motion}{moving around a circle at a constant speed.}
  \kw{Period ($T$)}{the time for one full turn. Unit: second (s).}
  \kw{Tangent}{a line that touches the circle at one point, at \ang{90} to the radius.}
  \kw{Centripetal acceleration ($a_\mathrm{c}$)}{the acceleration towards the centre of the circle.}
  \kw{Centripetal force ($F_\mathrm{c}$)}{the net force towards the centre of the circle.}
\end{keywords}

A car drives around a roundabout at a steady \qty{30}{km.h^{-1}}. The speedometer does not
change, but the direction keeps changing. This is \term{uniform circular motion}: motion
around a circle at a constant speed. ``Uniform'' means ``not changing''.

\begin{refresher}{Velocity, acceleration and Newton's laws}
Velocity is speed \emph{and} direction. It is a \term{vector}: a quantity with a size and a
direction. If the speed or the direction changes, the velocity changes. A change in velocity
is an acceleration. The \term{net force} is the total of all the forces on an object, added
with their directions. First law: with no net force, an object stays still or moves in a
straight line at a constant speed. Second law: $\vec{F}_\mathrm{net} = m\vec{a}$ (on the
sheet), so the net force and the acceleration point the same way. The \term{weight} of an
object is the pull of gravity on it: weight $= mg$, with $g = \qty{9.8}{m.s^{-2}}$ (data
sheet). A \term{free-body diagram} is a simple drawing of one object, with one arrow for each
force on it.
\end{refresher}

\begin{figure}
  \centering
  \begin{tikzpicture}[x=1cm, y=1cm]
    \begin{scope}
      \draw[hscmuted!55, line width=0.8pt] (0,0) circle (1.4);
      \fill[hscink] (0,0) circle (1.5pt);
      \node[diagramlabel, below, inner sep=3pt] at (0,0) {centre};
      \draw[->, hscmuted, line width=0.6pt] (212:1.75) arc (212:255:1.75);
      \node[diagramlabel, text=hscmuted, anchor=north east] at (232:1.85) {motion};
      \draw[guide] (35:1.4) -- ++(125:2.4);
      \node[diagramlabel, text=hscmuted, anchor=south, align=center] at ($(35:1.4)+(125:2.35)+(0.35,0.05)$)
        {path if the\\ string breaks};
      \foreach \a in {35, 155, 275} {
        \coordinate (P) at (\a:1.4);
        \draw[velocity] (P) -- ++(\a+90:1.0);
        \draw[accel] (P) -- ($(P)!0.5!(0,0)$);
        \fill[object] (P) circle (4pt);
      }
      \node[diagramlabel, text=hscvel, anchor=west] at ($(35:1.4)+(125:0.9)+(0.08,0)$) {$v$};
      \node[diagramlabel, text=hscacc, anchor=north west, inner sep=1pt] at ($(35:0.85)+(0.05,-0.02)$) {$a_\mathrm{c}$};
      \node[diagramlabel, text=hscvel] at ($(155:1.4)+(245:1.2)$) {$v$};
      \node[diagramlabel, text=hscacc, anchor=south, inner sep=2pt] at ($(155:0.85)+(0.05,0.05)$) {$a_\mathrm{c}$};
      \node[diagramlabel, text=hscvel, anchor=west] at ($(275:1.4)+(5:1.05)$) {$v$};
      \node[diagramlabel, text=hscacc, anchor=west, inner sep=2pt] at ($(275:0.95)+(0.02,-0.1)$) {$a_\mathrm{c}$};
      \node[font=\small\bfseries, text=hscink] at (-1.8,1.8) {(a)};
    \end{scope}
    \begin{scope}[xshift=6.4cm, yshift=-0.8cm]
      \coordinate (O) at (0,0);
      \coordinate (A) at (110:1.8);
      \coordinate (B) at (160:1.8);
      \draw[velocity] (O) -- (A) node[pos=0.55, right, inner sep=2pt, text=hscvel] {$v_1$ (earlier)};
      \draw[velocity] (O) -- (B) node[pos=0.45, below left, inner sep=2pt, text=hscvel] {$v_2$ (later)};
      \draw[accel] (A) -- (B) node[pos=0.5, above left, inner sep=1pt, text=hscacc] {$\Delta v$};
      \fill[hscink] (O) circle (1.3pt);
      \node[font=\small\bfseries, text=hscink] at (-2.1,2.5) {(b)};
    \end{scope}
  \end{tikzpicture}
  \caption{(a) Top view of a ball whirled on a string (the string is not drawn). The
    velocity $v$ is along the tangent. The acceleration $a_\mathrm{c}$ points to the centre.
    (b) The velocities at two nearby points, a moment apart, drawn tail to tail. The change
    $\Delta v$ points towards the centre of the circle.}
  \label{ch2:fig:ucm}
\end{figure}

Look at \cref{ch2:fig:ucm}a. At every moment, the velocity lies along a \term{tangent}. A
tangent is a line that touches the circle at one point. It is at \ang{90} to the radius.

The velocity arrow keeps turning. In \cref{ch2:fig:ucm}b, $v_1$ is the velocity at one
moment and $v_2$ is the velocity a short time later. Both arrows have the same length (the
same speed), but they point in different directions.

The change in velocity is $\Delta v = v_2 - v_1$ ($\Delta$ means ``change in''). To find it,
draw the two arrows tail to tail. Then $\Delta v$ is the arrow from the tip of $v_1$ to the
tip of $v_2$, because $v_1 + \Delta v = v_2$. It points towards the centre of the circle.

Acceleration is the change in velocity divided by the time taken. So the acceleration also
points to the centre. It is the \term{centripetal acceleration} ($a_\mathrm{c}$): the
acceleration towards the centre. ``Centripetal'' means ``centre-seeking''.

In one full turn, the object goes once around the circle. The distance once around a circle
is its \term{circumference}, $2\pi r$. The time for one full turn is the \term{period} ($T$).
Speed is distance divided by time, so $v = \dfrac{2\pi r}{T}$.

\begin{toolkit}{Units and powers of ten}
\si{km.h^{-1}} to \si{m.s^{-1}}: divide by 3.6 (\qty{72}{km.h^{-1}} $= \qty{20}{m.s^{-1}}$).
Diameter to radius: halve it. \si{cm} to \si{m}: divide by 100. \si{g} to \si{kg}: divide by
1000. \term{Frequency} ($f$) is the number of full turns each second. Its unit is the hertz
(\si{Hz}): \qty{1}{Hz} is one turn per second. Revolutions per minute (rpm) to $f$: divide
by 60. Then $T = \dfrac{1}{f}$ (from $f = \dfrac{1}{T}$ on the sheet). Example:
\qty{1200}{rpm} gives $f = 1200 \div 60 = \qty{20}{Hz}$, so $T = \dfrac{1}{20} = \qty{0.050}{s}$.

Answers are often written in \term{scientific notation}: a number from 1 to 10, times a power
of ten. Examples: $2.9 \times 10^2 = 2.9 \times 100 = 290$, and
$1.3 \times 10^4 = 1.3 \times 10\,000 = 13\,000$. This form also shows the significant
figures: $2.9 \times 10^2$ has 2.
\end{toolkit}

\begin{formulacard}{Speed in a circle}
  \cardformula{v = \frac{2\pi r}{T}}
  \begin{symbols}
    \sym{$v$}{speed of the object (along the tangent)}{\si{m.s^{-1}}}
    \sym{$r$}{radius of the circle}{\si{m}}
    \sym{$T$}{period: the time for one full turn}{\si{s}}
  \end{symbols}
  \inwords{speed = distance for one lap ($2\pi r$) divided by the time for one lap ($T$).}
  \whentouse{an object goes around a circle at a constant speed.}
  \whennottouse{the speed changes during the turn (it then gives only the average speed).}
  \rearranged{$T = \dfrac{2\pi r}{v}$ \qquad $r = \dfrac{vT}{2\pi}$ \qquad
    with $f = \dfrac{1}{T}$: \ $v = 2\pi r f$}
  \tensecond{$r = \qty{3.0}{m}$, $T = \qty{6.0}{s}$, so
    $v = \dfrac{2\pi \times 3.0}{6.0} = \qty{3.1}{m.s^{-1}}$.}
  \trap{here $T$ is the period (not tension). Use the radius, not the diameter, and the
    time for \emph{one} turn in seconds.}
\end{formulacard}

How big is $a_\mathrm{c}$? A faster object changes its velocity in two ways at once.
(1) It goes round sooner, so its velocity arrow turns more quickly. (2) Its arrow is longer,
so each bit of turning is a bigger change. Double the speed, and each effect doubles.
Together, $a_\mathrm{c}$ becomes $2 \times 2 = 4$ times bigger. So $a_\mathrm{c}$ grows with
$v^2$. A tighter circle (smaller $r$) also turns the arrow more quickly, so $a_\mathrm{c}$ is
bigger.

\begin{formulacard}{Centripetal acceleration}
  \cardformula{a_\mathrm{c} = \frac{v^2}{r}}
  \begin{symbols}
    \sym{$a_\mathrm{c}$}{centripetal acceleration, towards the centre}{\si{m.s^{-2}}}
    \sym{$v$}{speed of the object}{\si{m.s^{-1}}}
    \sym{$r$}{radius of the circle}{\si{m}}
  \end{symbols}
  \inwords{the acceleration towards the centre grows with the square of the speed and is
    bigger for a tighter circle.}
  \whentouse{any object moving in a circle (at that moment).}
  \whennottouse{to find how fast the speed changes: $a_\mathrm{c}$ changes only the direction.}
  \rearranged{$v = \sqrt{a_\mathrm{c}\, r}$ \qquad $r = \dfrac{v^2}{a_\mathrm{c}}$.
    With $v = \dfrac{2\pi r}{T}$: \
    $a_\mathrm{c} = \dfrac{(2\pi r/T)^2}{r} = \dfrac{4\pi^2 r^2}{T^2\, r} = \dfrac{4\pi^2 r}{T^2}$
    (not printed in this form: combine the two sheet formulas $v = \dfrac{2\pi r}{T}$ and
    $a_\mathrm{c} = \dfrac{v^2}{r}$)}
  \tensecond{$v = \qty{4.0}{m.s^{-1}}$, $r = \qty{2.0}{m}$, so
    $a_\mathrm{c} = \dfrac{4.0^2}{2.0} = \qty{8.0}{m.s^{-2}}$.}
  \trap{constant speed does \emph{not} mean zero acceleration. $a_\mathrm{c}$ points to
    the centre, never outwards.}
\end{formulacard}

The sheet also lists $\omega = \dfrac{\Delta\theta}{t}$. \term{Angular velocity} ($\omega$)
is the angle turned each second. Here angles are measured in \term{radians} (rad), another
unit for angles. One full turn is $2\pi$ rad, the same as \ang{360}. In one period the angle
turned is $2\pi$ rad. Put $\Delta\theta = 2\pi$ and $t = T$ into the sheet formula:
$\omega = \dfrac{2\pi}{T}$ (not on the sheet). Unit: \si{rad.s^{-1}}. \Cref{ch3:sec:omega}
explains $\omega$ in full.

By Newton's second law, an acceleration towards the centre needs a net force towards the
centre: $F_\mathrm{net} = ma_\mathrm{c} = \dfrac{mv^2}{r}$. This net force is the
\term{centripetal force} ($F_\mathrm{c}$). Centripetal force is not a new kind of force. It
is a job that real forces do: the job of pulling the object towards the centre.

Always ask: which real force does the job? \Cref{ch2:tab:jobs} gives the answers. Do not
worry about the new names in it yet: \cref{ch2:sec:situations} explains each situation with
a diagram.

\begin{table}
  \caption{Which real force provides the centripetal force?}\label{ch2:tab:jobs}
  \centering
  \begin{tabular}{@{}L{78mm}L{62mm}@{}}
    \toprule
    \colhead{Situation} & \colhead{Real force towards the centre}\\
    \midrule
    Ball on a string on a smooth (frictionless) table & tension (the pull of the string)\\
    Conical pendulum (a ball on a string going round in a horizontal circle);
      swing-ride chairs & horizontal part of the tension\\
    Playground swing, at its lowest point & tension minus weight ($F_\mathrm{T} - mg$)\\
    Car on a banked track (a road tilted up at the outside of the bend), at its design
      speed; marble in a bowl & horizontal part of the normal force (the push of a
      surface)\\
    Car on a flat bend & sideways friction from the road on the tyres\\
    Rider in a rotor ride (a spinning drum) & normal force from the wall\\
    Moon or satellite in orbit (\cref{ch4:chap}) & gravity\\
    \bottomrule
  \end{tabular}
\end{table}

\begin{formulacard}{Centripetal force}
  \cardformula{F_\mathrm{c} = \frac{mv^2}{r}}
  \begin{symbols}
    \sym{$F_\mathrm{c}$}{centripetal force: the net force towards the centre}{\si{N}}
    \sym{$m$}{mass of the object}{\si{kg}}
    \sym{$v$}{speed of the object}{\si{m.s^{-1}}}
    \sym{$r$}{radius of the circle}{\si{m}}
  \end{symbols}
  \inwords{the inward net force needed grows with the mass and with the square of the
    speed, and is bigger for a tighter circle.}
  \whentouse{any object moving in a circle ($F_\mathrm{net} = ma$ with $a_\mathrm{c} = \dfrac{v^2}{r}$).}
  \whennottouse{as an extra arrow on a free-body diagram, or for forces along the motion.}
  \rearranged{$v = \sqrt{\dfrac{F_\mathrm{c}\, r}{m}}$ \qquad $r = \dfrac{mv^2}{F_\mathrm{c}}$
    \qquad $m = \dfrac{F_\mathrm{c}\, r}{v^2}$}
  \tensecond{$m = \qty{0.50}{kg}$, $v = \qty{6.0}{m.s^{-1}}$, $r = \qty{3.0}{m}$, so
    $F_\mathrm{c} = \dfrac{0.50 \times 6.0^2}{3.0} = \qty{6.0}{N}$.}
  \trap{$F_\mathrm{c}$ is the \emph{net} force, made by real forces. Never label an arrow
    ``centripetal force''. Mass must be in \si{kg}, not \si{g}.}
\end{formulacard}

When a car turns left, you feel pushed to the right. But no force pushes you outwards.
Because of \term{inertia} (the tendency to keep moving in a straight line), your body tries
to keep going straight. The seat (by friction), the seatbelt and, if you slide across, the
door push you \emph{inwards}. So you turn with the car.

Some people call the outward feeling a ``centrifugal force'' (``centrifugal'' means
``centre-fleeing''). It is not a real force. Do not use it in this course.

If a whirling ball's string breaks, the inward force is gone. By Newton's first law, the
ball then moves in a straight line along the tangent. It does not fly straight outwards
(\cref{ch2:fig:ucm}a).

\begin{toolkit}{Proportional reasoning: ``what if \ldots?''}
The sign $\propto$ means ``is proportional to''. $F_\mathrm{c} \propto m$ means: double $m$
and $F_\mathrm{c}$ doubles; triple $m$ and $F_\mathrm{c}$ triples.
In $F_\mathrm{c} = \dfrac{mv^2}{r}$, change one quantity and keep the others the same.
Mass $\times 2$ gives $F_\mathrm{c} \times 2$ ($F_\mathrm{c} \propto m$). Speed $\times 2$ gives
$F_\mathrm{c} \times 2^2 = \times 4$ ($F_\mathrm{c} \propto v^2$). Radius $\div 2$ at the same
speed gives $F_\mathrm{c} \times 2$ ($F_\mathrm{c} \propto \dfrac{1}{r}$). You test these
relationships in the investigation (\cref{ch2:sec:exp}).
\end{toolkit}

\begin{workedexample}{A cyclist on a circular track}\label{ch2:we:cyclist}
A cyclist rides around a flat (not tilted) circular track of radius \qty{32}{m} at a
constant speed. Each lap
takes \qty{18}{s}. The cyclist and bike have a total mass of \qty{75}{kg}. Calculate (a) the
speed, (b) the centripetal acceleration and (c) the net force on the cyclist and bike.
\given{$r = \qty{32}{m}$, \ $T = \qty{18}{s}$ (one lap = one period), \ $m = \qty{75}{kg}$}
\find{(a) $v$ \quad (b) $a_\mathrm{c}$ \quad (c) $F_\mathrm{c}$}
\formula{$v = \dfrac{2\pi r}{T}$, \quad $a_\mathrm{c} = \dfrac{v^2}{r}$, \quad
  $F_\mathrm{c} = ma_\mathrm{c}$}
\substitute{\eqlines{v &= \frac{2\pi \times \qty{32}{m}}{\qty{18}{s}} = \qty{11.17}{m.s^{-1}}
    &&\text{(keep the extra digits)}\\
  a_\mathrm{c} &= \frac{(\qty{11.17}{m.s^{-1}})^2}{\qty{32}{m}} = \qty{3.899}{m.s^{-2}}\\
  F_\mathrm{c} &= \qty{75}{kg} \times \qty{3.899}{m.s^{-2}} = \qty{292}{N}}}
\answer{(a) \final{$v = \qty{11}{m.s^{-1}}$} \ (b) \final{$a_\mathrm{c} = \qty{3.9}{m.s^{-2}}$}
  \ (c) \final{$F_\mathrm{c} = \qty{2.9e2}{N}$} towards the centre}
\checkstep{\qty{11}{m.s^{-1}} is about \qty{40}{km.h^{-1}}, a fast cyclist. The
  \qty{290}{N} comes from sideways friction on the tyres (\cref{ch2:sec:situations}).}
\end{workedexample}

\begin{workedexample}{A puck on an air table}\label{ch2:we:puck}
A \qty{0.40}{kg} puck on a frictionless air table is tied to a fixed peg by a string. It
moves in a circle of radius \qty{0.75}{m} at \qty{2.5}{m.s^{-1}}. The string breaks if the
tension is more than \qty{5.0}{N}. (a) Find the tension. (b) The string is shortened to
\qty{0.375}{m} at the same speed. Does it break? (c) Describe the path after it breaks.
\given{$m = \qty{0.40}{kg}$, \ $r = \qty{0.75}{m}$, \ $v = \qty{2.5}{m.s^{-1}}$, \
  breaking tension \qty{5.0}{N}. The tension is the only horizontal force.}
\find{(a) $F_\mathrm{T}$ \quad (b) the new tension \quad (c) the path}
\formula{$F_\mathrm{T} = F_\mathrm{c} = \dfrac{mv^2}{r}$}
\substitute{\eqlines{\text{(a)}\quad F_\mathrm{T} &= \frac{\qty{0.40}{kg} \times (\qty{2.5}{m.s^{-1}})^2}{\qty{0.75}{m}}
    = \qty{3.33}{N}\\
  \text{(b)}\quad F_\mathrm{T} &= \frac{\qty{0.40}{kg} \times (\qty{2.5}{m.s^{-1}})^2}{\qty{0.375}{m}}
    = \qty{6.67}{N}\\
  &\text{(radius halved, so force doubled)}}}
\answer{(a) \final{$F_\mathrm{T} = \qty{3.3}{N}$} \ (b) \final{Yes: $\qty{6.7}{N} > \qty{5.0}{N}$}
  \ (c) a straight line along the tangent at \qty{2.5}{m.s^{-1}}.}
\checkstep{(b) fits the ``what if'' rule: $r \div 2$ gives $F_\mathrm{c} \times 2$.
  (c) With no friction there is no net force, so the speed stays the same.}
\end{workedexample}

\begin{mistakes}
\begin{itemize}
  \item \xmark\ ``Constant speed, so no acceleration.'' \cmark\ The direction changes: it accelerates.
  \item \xmark\ Drawing an ``$F_\mathrm{c}$'' or ``centrifugal'' arrow. \cmark\ Draw only real forces.
  \item \xmark\ ``If the string breaks, it flies straight out.'' \cmark\ It moves along the tangent.
  \item \xmark\ ``Double the speed, double the force.'' \cmark\ $v$ is squared: 4 times.
\end{itemize}
\end{mistakes}

\begin{examtip}
``Explain why an object moving in a circle at a constant speed is accelerating''
(2--3 marks). Full marks need three linked points:
\begin{enumerate}
  \item Velocity is a vector: it has a size and a direction.
  \item The direction keeps changing, so the velocity changes. That is an acceleration.
  \item The acceleration points to the centre. Its size is $\dfrac{v^2}{r}$.
\end{enumerate}
``Identify the force that provides the centripetal force'': name the \emph{real} force (for
example ``friction from the road on the tyres'') and say it acts towards the centre. The
answers ``centripetal force'' or ``centrifugal force'' earn no marks.
\end{examtip}

\begin{quickcheck}
  \item A Ferris wheel of radius \qty{20}{m} turns once every \qty{80}{s}. Find the speed
        of a passenger and the wheel's angular velocity $\omega = \dfrac{2\pi}{T}$.
  \item A car drives around a roundabout at a constant \qty{25}{km.h^{-1}}. Is it
        accelerating? Explain, and give the direction.
  \item A washing machine drum of diameter \qty{48}{cm} spins at \qty{1200}{rpm}. Find the
        period, the speed of a sock on the wall and its centripetal acceleration. Compare
        the acceleration with $g = \qty{9.8}{m.s^{-2}}$.
  \item A car needs \qty{3.0}{kN} of friction to take a bend at \qty{50}{km.h^{-1}}. How
        much does it need (a) at \qty{100}{km.h^{-1}} on the same bend; (b) at
        \qty{50}{km.h^{-1}} on a bend with half the radius?
  \item Name the real force that provides the centripetal force: (a) the Moon orbiting
        the Earth; (b) clothes in a spinning drum; (c) a car on a flat roundabout.
\end{quickcheck}

%% ==================================================================
\section{What provides the centripetal force? Real situations}\label{ch2:sec:situations}
\covers{M5.IQ2.02, M5.IQ2.03}

\begin{bigidea}
In every situation, find the real forces and split them into a part towards the centre
and a vertical part. The part towards the centre adds up to $mv^2/r$. Strings pull along
their length, surfaces push at \ang{90} to themselves, and friction acts along a surface.
\end{bigidea}

\begin{keywords}
  \kw{Tension ($F_\mathrm{T}$)}{the pull of a string, rope or chain, along its length.}
  \kw{Normal force ($F_\mathrm{N}$)}{the push of a surface on an object, at \ang{90} to the surface.}
  \kw{Component}{the part of a force that acts in one chosen direction.}
  \kw{Conical pendulum}{a mass hanging on a string from a fixed point and moving in a
    horizontal circle below it; the string sweeps out a cone.}
  \kw{Banked track}{a curved road that is tilted, higher on the outside of the bend.}
  \kw{Static friction ($F_\mathrm{f}$)}{friction that stops sliding; it grows as needed, up
    to a maximum.}
  \kw{Rotor}{a ride: a spinning drum whose riders stay on the wall when the floor drops.}
\end{keywords}

\textbf{A method for every circular-motion force problem.}
\begin{enumerate}
  \item Find the circle: where is its centre, and what is $r$?
  \item Draw a free-body diagram: one arrow for each real force (weight, tension, normal
        force, friction). Do not draw an $F_\mathrm{c}$ arrow.
  \item Split each force into a part towards the centre and a vertical part.
  \item Write two equations. Towards the centre: total $= \dfrac{mv^2}{r}$. Vertically (for a
        horizontal circle): total $= 0$. Then solve.
\end{enumerate}

\begin{toolkit}{Splitting a force into components}
\begin{sidebyside}[0.72]
A force $F$ makes an angle $\theta$ with the \emph{vertical}. Its vertical part is
$F\cos\theta$: this side is next to the angle, so use $\cos$. Its horizontal part is
$F\sin\theta$: this side is opposite the angle, so use $\sin$. Check: at $\theta = \ang{0}$
the force is vertical, and $F\cos 0^\circ = F$ \cmark. If $\theta$ is measured from the
\emph{horizontal}, swap $\sin$ and $\cos$. Since $\dfrac{\sin\theta}{\cos\theta} = \tan\theta$,
dividing the parts gives $\dfrac{F\sin\theta}{F\cos\theta} = \tan\theta$. To find an angle
from its tan, use the $\tan^{-1}$ key, with the calculator in degree mode.
\nextside
  \centering
  \begin{tikzpicture}[x=1cm, y=1cm]
    \draw[guide] (0,0) -- (0,1.6);
    \draw[guide] (0,1.6) -- (1.12,1.6);
    \draw[force] (0,0) -- (1.12,1.6) node[pos=0.5, right, text=hscforce] {$F$};
    \draw[hscink, line width=0.5pt] (0,0.55) arc (90:55:0.55);
    \node[diagramlabel] at (0.18,0.8) {$\theta$};
    \node[diagramlabel, anchor=east] at (0,0.9) {$F\cos\theta$};
    \node[diagramlabel, anchor=south] at (0.56,1.6) {$F\sin\theta$};
  \end{tikzpicture}
\end{sidebyside}
\end{toolkit}

\subsection*{Strings: conical pendulums and swings}
Hang a ball on a string and push it so that it goes round in a horizontal circle. The
string sweeps out a cone: this is a \term{conical pendulum}. A ceiling light swinging in a
circle moves like this. Only two real forces act on the ball (\cref{ch2:fig:conical}): the
\term{tension} $F_\mathrm{T}$ (the string's pull) and the weight $mg$.

The string hangs from a fixed point, the \term{pivot}. The string has length $L$ and makes an
angle $\theta$ with the vertical. The centre of the circle is directly \emph{below} the
pivot, not at it. In \cref{ch2:fig:conical}a, the string is the long side (the hypotenuse) of
a right-angled triangle, and $r$ is the side opposite $\theta$. Opposite $=$ hypotenuse
$\times \sin\theta$, so $r = L\sin\theta$.

\begin{figure}
  \centering
  \begin{tikzpicture}[x=1cm, y=1cm]
    \begin{scope}
      \fill[ground] (-1.5,0) rectangle (2.4,0.22);
      \draw[hscink, line width=0.8pt] (-1.5,0) -- (2.4,0);
      \coordinate (piv) at (0,0);
      \coordinate (B) at (1.548,-2.212);
      \coordinate (C) at (0,-2.212);
      \draw[guide] (piv) -- (0,-2.7);
      \draw[hscmuted, dashed, line width=0.7pt] (C) ellipse (1.548 and 0.36);
      \draw[hscink, line width=1pt] (piv) -- (B) node[pos=0.5, right, diagramlabel] {$L$};
      \draw[hscink, line width=0.5pt] (0,-0.7) arc (270:305:0.7);
      \node[diagramlabel] at (0.19,-0.92) {$\theta$};
      \draw[displacement, <->] (C) -- (B) node[pos=0.5, below, diagramlabel, inner sep=2pt] {$r$};
      \fill[hscink] (C) circle (1.5pt);
      \draw[hscmuted, line width=0.4pt] (C) -- (-1.75,-2.75);
      \node[diagramlabel, anchor=north east, align=right, inner sep=1pt] at (-1.75,-2.75) {centre of\\the circle};
      \fill[object] (B) circle (5pt);
      \fill[hscink] (piv) circle (1.5pt);
      \node[diagramlabel, anchor=east] at (-0.12,-0.25) {pivot};
      \node[font=\small\bfseries, text=hscink] at (-1.3,-0.75) {(a)};
    \end{scope}
    \begin{scope}[xshift=7.4cm, yshift=-1.7cm]
      \coordinate (O) at (0,0);
      \draw[guide] (O) -- (0,1.5);
      \draw[guide] (0,1.5) -- (-1.05,1.5);
      \draw[force] (O) -- (-1.05,1.5) node[pos=0.62, below left, inner sep=1pt, text=hscforce] {$F_\mathrm{T}$};
      \draw[force] (O) -- (0,-1.5) node[pos=0.8, right, text=hscforce] {$mg$};
      \draw[hscink, line width=0.5pt] (0,0.55) arc (90:125:0.55);
      \node[diagramlabel] at (-0.15,0.8) {$\theta$};
      \node[diagramlabel, anchor=west] at (0.05,1.05) {$F_\mathrm{T}\cos\theta$};
      \node[diagramlabel, anchor=south] at (-0.53,1.53) {$F_\mathrm{T}\sin\theta$};
      \draw[displacement] (-1.4,-0.6) -- (-2.5,-0.6);
      \node[diagramlabel, anchor=north, text=hscmuted] at (-1.95,-0.65) {to the centre};
      \fill[object] (O) circle (5pt);
      \node[font=\small\bfseries, text=hscink] at (-2.5,1.8) {(b)};
    \end{scope}
  \end{tikzpicture}
  \caption{A conical pendulum. (a) Side view: the centre of the circle is directly below
    the pivot. (b) Free-body diagram: only tension and weight act. Dashed lines are the
    components of the tension.}
  \label{ch2:fig:conical}
\end{figure}

Split the tension into \term{components} (parts in chosen directions), with $\theta$ from the
vertical (see the Maths toolkit above). The weight has no horizontal part. So the horizontal
part of the tension is the whole net force towards the centre:
$F_\mathrm{T}\sin\theta = \dfrac{mv^2}{r}$. The ball does not move up or down, so the vertical
forces balance (up $=$ down): $F_\mathrm{T}\cos\theta = mg$.

The chairs of a swing ride hang in the same way. But their chains hang from a spinning
frame, not from the axis. So for a swing-ride chair, $r = d + L\sin\theta$, where $d$ is the
distance from the axis to the top of the chain.

A \term{vertical circle} is an upright circle, like the path of a playground swing. The
speed changes as the object goes up and down. So this is \emph{not} uniform circular motion.
But at any one point, the net force towards the centre is still $\dfrac{mv^2}{r}$.

At the lowest point of a swing, the centre of the circle is \emph{above} the child. Take
``towards the centre'' (up) as positive. The tension points up, to the centre: $+F_\mathrm{T}$.
The weight points down, away from the centre: $-mg$. So $F_\mathrm{T} - mg = \dfrac{mv^2}{r}$.
The chains pull harder than your weight, so you feel heavier.

At the top of a vertical circle on a string, both forces point down, towards the centre:
$F_\mathrm{T} + mg = \dfrac{mv^2}{r}$. These two equations are not on the formula sheet. Each
one comes from ``net force towards the centre $= \dfrac{mv^2}{r}$'' at that point.
\Cref{ch3:sec:work} covers vertical circles in full.

\begin{workedexample}{A child on a swing}\label{ch2:we:swing}
A \qty{25}{kg} child sits on a swing with chains \qty{2.0}{m} long. At the lowest point the
child moves at \qty{3.0}{m.s^{-1}}. Find the total tension in the chains at that point.
\given{$m = \qty{25}{kg}$, \ $v = \qty{3.0}{m.s^{-1}}$, \ $r = \qty{2.0}{m}$ (treat the child
  as a point at the end of the chains)}
\find{the total tension $F_\mathrm{T}$ at the lowest point}
\formula{towards the centre (up): $F_\mathrm{T} - mg = \dfrac{mv^2}{r}$. Add $mg$ to both
  sides: $F_\mathrm{T} = mg + \dfrac{mv^2}{r}$}
\substitute{\eqlines{mg &= \qty{25}{kg} \times \qty{9.8}{m.s^{-2}} = \qty{245}{N}\\
  \frac{mv^2}{r} &= \frac{\qty{25}{kg} \times (\qty{3.0}{m.s^{-1}})^2}{\qty{2.0}{m}} = \qty{112.5}{N}\\
  F_\mathrm{T} &= \qty{245}{N} + \qty{112.5}{N} = \qty{357.5}{N}}}
\answer{\final{$F_\mathrm{T} = \qty{3.6e2}{N}$}, upwards (2 significant figures)}
\checkstep{$F_\mathrm{T}$ is bigger than the weight: $\dfrac{357.5}{245} = 1.46$, so the child
  feels about 1.5 times as heavy \cmark. If the swing were still ($v = 0$), $F_\mathrm{T}$
  would equal $mg$ \cmark.}
\end{workedexample}

\subsection*{Tilted surfaces: bowls and banked tracks}
Roll a marble around the inside of a bowl. No string holds it. The wall pushes on it
instead. This push is the \term{normal force} $F_\mathrm{N}$: the push of a surface, always at
\ang{90} to the surface. On the sloping wall of the bowl, $F_\mathrm{N}$ tilts inwards,
towards the centre of the circle (\cref{ch2:fig:bank}a).

In a round bowl, $F_\mathrm{N}$ lies along the line from the marble to the bowl's centre C.
So the angle $\theta$ marked at C is the same as the angle of $F_\mathrm{N}$ from the
vertical. The radius $r$ of the marble's circle is its horizontal distance from the centre
line. It is not the radius of the bowl.

A \term{banked track} (a bend that is higher on the outside) works in the same way for a car
(\cref{ch2:fig:bank}b).

\begin{figure}
  \centering
  \begin{tikzpicture}[x=1cm, y=1cm]
    \begin{scope}
      \coordinate (C) at (0,0);
      \draw[hscink, line width=1.6pt] (180:1.8) arc (180:360:1.8);
      \draw[guide] (0,0.3) -- (0,-2.0);
      \coordinate (M) at (1.379,-1.157);
      \draw[guide] (C) -- (M);
      \draw[hscink, line width=0.5pt] (0,-0.45) arc (270:320:0.45);
      \node[diagramlabel] at (290:0.66) {$\theta$};
      \draw[displacement, <->] (0,-1.157) -- (M) node[pos=0.4, below, diagramlabel] {$r$};
      % F_N along M -> C (140 deg); length 1.15, so its vertical part 1.15 sin140 = 0.74 = mg arrow
      \draw[force] (M) -- ++(140:1.15) node[pos=0.75, above right, inner sep=1pt, text=hscforce] {$F_\mathrm{N}$};
      \draw[force] (M) -- ++(0,-0.74) node[pos=0.85, right, text=hscforce] {$mg$};
      \fill[object] (M) circle (4.5pt);
      \fill[hscink] (C) circle (1.5pt);
      \node[diagramlabel, anchor=south east] at (C) {C};
      \node[font=\small\bfseries, text=hscink] at (-1.9,0.5) {(a)};
    \end{scope}
    \begin{scope}[xshift=5.4cm, yshift=-1.8cm]
      \fill[ground] (0,0) -- (4.2,1.529) -- (4.2,1.2) -- (0,-0.33) -- cycle;
      \draw[hscink, line width=1pt] (0,0) -- (4.2,1.529);
      \draw[guide] (0,0) -- (3.0,0);
      \draw[hscink, line width=0.5pt] (1.0,0) arc (0:20:1.0);
      \node[diagramlabel] at (1.25,0.17) {$\theta$};
      \coordinate (P) at (2.632,0.958);
      \begin{scope}[rotate around={20:(P)}]
        \draw[object, rounded corners=1.5pt] ($(P)+(-0.55,0.12)$) rectangle ($(P)+(0.55,0.62)$);
        \fill[hscink] ($(P)+(-0.33,0.11)$) circle (0.11);
        \fill[hscink] ($(P)+(0.33,0.11)$) circle (0.11);
      \end{scope}
      \coordinate (G) at ($(P)+(110:0.37)$);
      \draw[guide] (G) -- ++(0,1.6);
      \draw[guide] ($(G)+(0,1.6)$) -- ++(-0.582,0);
      \draw[force] (G) -- ++(110:1.7) node[pos=0.95, left, inner sep=2pt, text=hscforce] {$F_\mathrm{N}$};
      \draw[force] (G) -- ++(0,-1.6) node[pos=0.9, right, text=hscforce] {$mg$};
      \draw[hscink, line width=0.5pt] ($(G)+(0,0.7)$) arc (90:110:0.7);
      \node[diagramlabel] at ($(G)+(-0.1,0.92)$) {$\theta$};
      \draw[displacement] (1.3,1.3) -- (0.1,1.3);
      \node[diagramlabel, anchor=south, text=hscmuted, align=center] at (0.7,1.35) {to the centre\\of the bend};
      \node[font=\small\bfseries, text=hscink] at (-0.1,2.85) {(b)};
    \end{scope}
  \end{tikzpicture}
  \caption{(a) A marble circling inside a round bowl (cross-section): $F_\mathrm{N}$ points to
    the bowl's centre C. (b) A car on a banked track, seen from behind: $F_\mathrm{N}$ is at
    \ang{90} to the road. In both, $\theta$ is the angle of $F_\mathrm{N}$ from the vertical.}
  \label{ch2:fig:bank}
\end{figure}

The maths is the same as for the string, with $F_\mathrm{N}$ in place of $F_\mathrm{T}$.
Towards the centre: $F_\mathrm{N}\sin\theta = \dfrac{mv^2}{r}$. Vertically:
$F_\mathrm{N}\cos\theta = mg$.

On a road, $\theta$ is also the \term{bank angle}: the road's tilt from the horizontal. Why?
$F_\mathrm{N}$ is at \ang{90} to the road, and the vertical is at \ang{90} to the horizontal.
Tilt the road by $\theta$, and $F_\mathrm{N}$ tilts by the same $\theta$ from the vertical
(both angles are marked in \cref{ch2:fig:bank}b).

Divide both sides of $F_\mathrm{N}\cos\theta = mg$ by $\cos\theta$:
$F_\mathrm{N} = \dfrac{mg}{\cos\theta}$. Here $\cos\theta$ is less than 1, so $F_\mathrm{N}$ is
\emph{bigger} than the weight.

At one speed, the \term{design speed}, no friction is needed. If the car goes faster, it
tends to slide up the slope, so friction acts \emph{down} the slope. If the car goes slower,
it tends to slide down, so friction acts \emph{up} the slope. Banking means less friction is
needed, so turns in wet weather are safer. \Cref{ch2:fig:bankfriction} shows the three cases.

\begin{figure}
  \centering
  % Added by the coverage audit (M5.IQ2.02: banked track at other speeds had no diagram).
  % Bank angle 20 deg; road (0,0)--(2.8,1.019); car at P = (1.4,0.510);
  \begin{tikzpicture}[x=1.2cm, y=1.2cm]
    \foreach \dx in {0, 4.0, 8.0} {
      \begin{scope}[xshift=\dx*1.2cm]
        \fill[ground] (0,0) -- (2.8,1.019) -- (2.8,0.70) -- (0,-0.319) -- cycle;
        \draw[hscink, line width=1pt] (0,0) -- (2.8,1.019);
        \draw[guide] (0,0) -- (0.9,0);
        \draw[hscink, line width=0.5pt] (0.5,0) arc (0:20:0.5);
        \node[diagramlabel, font=\footnotesize] at (0.66,0.1) {$\theta$};
        \begin{scope}[rotate around={20:(1.4,0.510)}]
          \draw[object, rounded corners=1.5pt] (0.98,0.60) rectangle (1.82,0.98);
          \fill[hscink] (1.15,0.60) circle (0.085);
          \fill[hscink] (1.65,0.60) circle (0.085);
        \end{scope}
        \draw[force] (1.314,0.744) -- ++(0,-1.0) node[pos=0.9, right, text=hscforce, inner sep=2pt] {$mg$};
        \draw[force] (1.314,0.744) -- ++(110:1.064) node[pos=1, left, text=hscforce] {$F_\mathrm{N}$};
        \fill[hscink] (1.314,0.744) circle (1.3pt);
      \end{scope}
    }
    % friction: (a) up the slope, (c) down the slope
    \draw[force] (1.314,0.744) -- ++(20:0.75) node[pos=1, above, text=hscforce] {$F_\mathrm{f}$};
    \draw[force] (9.314,0.744) -- ++(200:0.75) node[pos=1, above left, text=hscforce, inner sep=1pt] {$F_\mathrm{f}$};
    \node[diagramlabel, font=\footnotesize, align=center, anchor=north] at (1.4,-0.5)
      {\textbf{(a) slower}\\friction up the slope};
    \node[diagramlabel, font=\footnotesize, align=center, anchor=north] at (5.4,-0.5)
      {\textbf{(b) at the design speed}\\no friction needed};
    \node[diagramlabel, font=\footnotesize, align=center, anchor=north] at (9.4,-0.5)
      {\textbf{(c) faster}\\friction down the slope};
  \end{tikzpicture}
  \caption{Forces on a car on a banked track, seen from behind. The centre of the bend is to the
    left. (a) Slower than the design speed. (b) At the design speed. (c) Faster than the design
    speed. Arrows are not to scale.}\label{ch2:fig:bankfriction}
\end{figure}

\begin{formulacard}[notonsheet]{Tilted string or surface}
  \cardformula{\tan\theta = \frac{v^2}{rg}}
  \begin{symbols}
    \sym{$\theta$}{angle of the string (or of $F_\mathrm{N}$) from the vertical; for a road,
      the bank angle}{\si{\degree} (calculator in DEG mode)}
    \sym{$v$}{speed (for a banked road: the design speed)}{\si{m.s^{-1}}}
    \sym{$r$}{radius of the horizontal circle}{\si{m}}
    \sym{$g$}{gravitational acceleration, \qty{9.8}{m.s^{-2}}}{\si{m.s^{-2}}}
  \end{symbols}
  \inwords{a faster speed or a tighter circle needs a bigger lean (a steeper bank).}
  \whentouse{a conical pendulum; a banked road at its design speed; an object circling
    inside a bowl or cone without friction.}
  \whennottouse{a banked road at other speeds (friction also acts), or a flat road.}
  \rearranged{$v = \sqrt{rg\tan\theta}$ \quad $F_\mathrm{T}$ or $F_\mathrm{N} = \dfrac{mg}{\cos\theta}$
    \quad conical pendulum: $r = L\sin\theta$; \ swing-ride chair: $r = d + L\sin\theta$
    ($d$ = axis to top of chain)}
  \tensecond{a bend of $r = \qty{50}{m}$ banked at \ang{20}:
    $v = \sqrt{50 \times 9.8 \times \tan 20^\circ} = \qty{13.4}{m.s^{-1}}$
    ($13.4 \times 3.6 = \qty{48}{km.h^{-1}}$).}
  \derive{use the two component equations from the free-body diagram ($F$ is $F_\mathrm{T}$
    or $F_\mathrm{N}$).
    \begin{align*}
      F\sin\theta &= \frac{mv^2}{r} && \text{horizontal: net force towards the centre}\\
      F\cos\theta &= mg && \text{vertical: forces balance}\\
      \frac{F\sin\theta}{F\cos\theta} &= \frac{mv^2/r}{mg} && \text{divide the first line by the second}\\
      \tan\theta &= \frac{v^2}{rg} && \text{$F$ cancels ($\sin\theta \div \cos\theta = \tan\theta$); $m$ cancels}
    \end{align*}}
  \trap{for a string, $r$ is not the string length: $r = L\sin\theta$. On a slope, do not
    split forces along the slope (Year 11 habit): the acceleration is horizontal. If a
    question gives a wall's angle from the \emph{vertical} (for example a cone with walls at
    \ang{20} to the vertical), $F_\mathrm{N}$ is at \ang{20} from the \emph{horizontal}, so here
    $\theta = \ang{90} - \ang{20} = \ang{70}$. Mark the angle on your diagram first.}
\end{formulacard}

\begin{workedexample}{A conical pendulum}\label{ch2:we:conical}
A \qty{0.30}{kg} ball hangs from a \qty{0.90}{m} string and moves in a horizontal circle.
The string makes \ang{35} with the vertical. Find (a) the tension, (b) the radius of the
circle, (c) the speed of the ball and (d) the period.
\given{$m = \qty{0.30}{kg}$, \ $L = \qty{0.90}{m}$, \ $\theta = \ang{35}$ from the vertical}
\find{$F_\mathrm{T}$, \ $r$, \ $v$ and $T$}
\formula{$F_\mathrm{T}\cos\theta = mg$, \ $r = L\sin\theta$, \ $\tan\theta = \dfrac{v^2}{rg}$,
  \ $v = \dfrac{2\pi r}{T}$}
\substitute{\eqlines{F_\mathrm{T} &= \frac{mg}{\cos\theta}
    \quad\text{(divide both sides of $F_\mathrm{T}\cos\theta = mg$ by $\cos\theta$)}\\
  &= \frac{\qty{0.30}{kg} \times \qty{9.8}{m.s^{-2}}}{\cos 35^\circ}
    = \frac{\qty{2.94}{N}}{0.8192} = \qty{3.589}{N}\\
  r &= L\sin\theta = \qty{0.90}{m} \times 0.5736 = \qty{0.5162}{m}}\par
  % layout QA: two aligned blocks, so the step can break across the page (it left a third of
  % page empty); the phantom keeps the = signs of both blocks in line
  \eqlines{\hphantom{F_\mathrm{T}}\llap{$v^2$} &= rg\tan\theta
    \quad\text{(multiply both sides of $\tan\theta = \tfrac{v^2}{rg}$ by $rg$)}\\
  v &= \sqrt{\qty{0.5162}{m} \times \qty{9.8}{m.s^{-2}} \times 0.7002} = \qty{1.882}{m.s^{-1}}
    \quad\text{(square root)}\\
  T &= \frac{2\pi r}{v}
    \quad\text{(rearrange $v = \tfrac{2\pi r}{T}$: multiply by $T$, divide by $v$)}\\
  &= \frac{2\pi \times \qty{0.5162}{m}}{\qty{1.882}{m.s^{-1}}} = \qty{1.723}{s}}}
\answer{(a) \final{$F_\mathrm{T} = \qty{3.6}{N}$} \ (b) \final{$r = \qty{0.52}{m}$}
  \ (c) \final{$v = \qty{1.9}{m.s^{-1}}$} \ (d) \final{$T = \qty{1.7}{s}$}}
\checkstep{$F_\mathrm{T}$ is bigger than the weight (\qty{2.94}{N}), because the string must
  also pull inwards. $r$ is shorter than $L$ \cmark.}
\end{workedexample}

\begin{workedexample}{Banking a test-track bend}\label{ch2:we:bankroad}
Engineers want cars to take a test-track bend of radius \qty{120}{m} at \qty{72}{km.h^{-1}}
with no friction needed. (a) Find the bank angle. (b) Find the normal force on a
\qty{1300}{kg} car at this speed.
\given{$r = \qty{120}{m}$, \ $v = 72 \div 3.6 = \qty{20}{m.s^{-1}}$, \ $m = \qty{1300}{kg}$}
\find{(a) $\theta$ \quad (b) $F_\mathrm{N}$}
\formula{$\tan\theta = \dfrac{v^2}{rg}$ \quad and \quad $F_\mathrm{N} = \dfrac{mg}{\cos\theta}$}
\substitute{\eqlines{\tan\theta &= \frac{(\qty{20}{m.s^{-1}})^2}{\qty{120}{m} \times \qty{9.8}{m.s^{-2}}}
    = \frac{400}{1176} = 0.3401\\
  \theta &= \tan^{-1}(0.3401) = 18.79^\circ \quad\text{(inverse tan; keep the extra digits)}\\
  F_\mathrm{N} &= \frac{\qty{1300}{kg} \times \qty{9.8}{m.s^{-2}}}{\cos 18.79^\circ}
    = \frac{\qty{12740}{N}}{0.9467} = \qty{1.346e4}{N}}}
\answer{(a) \final{$\theta = \ang{19}$} \quad (b) \final{$F_\mathrm{N} = \qty{1.3e4}{N}$},
  at \ang{90} to the road}
\checkstep{$F_\mathrm{N} = \qty{1.35e4}{N}$ is about 6\,\% bigger than the weight,
  $mg = \qty{1.27e4}{N}$, because $\dfrac{1}{\cos 18.79^\circ} = 1.06$ \cmark.
  $\tan\theta$ has no unit: \si{m^2.s^{-2}} over \si{m^2.s^{-2}} \cmark.}
\end{workedexample}

\subsection*{Friction: flat bends and the rotor ride}

\begin{refresher}{Friction}
\term{Static friction} stops surfaces sliding over each other. It is only as big as it needs
to be, up to a maximum: $F_\mathrm{f} \le \mu F_\mathrm{N}$ ($\le$ means ``less than or equal
to''). The \term{coefficient of friction} ($\mu$) is a number with no unit that says how
grippy two surfaces are. It is about 0.8 for tyres on a dry road and about 0.1 on ice. On a
flat road $F_\mathrm{N} = mg$, so the largest friction is $\mu mg$. The formula sheet does not
print $F_\mathrm{f} \le \mu F_\mathrm{N}$ (Year 11), so learn it. Questions give the value of
$\mu$.
\end{refresher}

\begin{sidebyside}[0.50]
  \centering
  \begin{tikzpicture}[x=1cm, y=1cm]
    \begin{scope}
      \fill[hscmuted!12] (15:1.7) arc (15:115:1.7) -- (115:2.4) arc (115:15:2.4) -- cycle;
      \draw[hscmuted, line width=0.7pt] (15:1.7) arc (15:115:1.7);
      \draw[hscmuted, line width=0.7pt] (15:2.4) arc (15:115:2.4);
      \fill[hscink] (0,0) circle (1.5pt);
      \node[diagramlabel, anchor=north] at (0,-0.05) {centre};
      \coordinate (Q) at (62:2.05);
      \draw[guide] (0,0) -- (Q);
      \draw[velocity] (Q) -- ++(152:1.0) node[above, text=hscvel] {$v$};
      \draw[force] (Q) -- ($(Q)!0.45!(0,0)$) node[pos=0.65, right, text=hscforce] {$F_\mathrm{f}$};
      \draw[object, rotate around={152:(Q)}] ($(Q)+(-0.27,-0.15)$) rectangle ($(Q)+(0.27,0.15)$);
      \node[font=\small\bfseries, text=hscink] at (-1.5,2.3) {(a)};
    \end{scope}
    \begin{scope}[xshift=3.7cm, yshift=1.0cm]
      \coordinate (O) at (0,0);
      \draw[force] (O) -- (0,1.1) node[above, text=hscforce] {$F_\mathrm{N}$};
      \draw[force] (O) -- (0,-1.1) node[below, text=hscforce] {$mg$};
      \draw[force] (O) -- (-1.0,0) node[left, text=hscforce] {$F_\mathrm{f}$};
      \fill[object] (O) circle (5pt);
      \node[font=\small\bfseries, text=hscink] at (-0.9,1.3) {(b)};
    \end{scope}
  \end{tikzpicture}
  \captionof{figure}{A car on a flat bend. (a) Top view. (b) Free-body diagram, seen from
    behind (centre to the left).}\label{ch2:fig:flatbend}
\nextside
  When you steer, the tyres push sideways on the road, outwards. Newton's third law says: if
  A pushes on B, then B pushes back on A with an equal force in the opposite direction. So
  the road pushes the tyres sideways, \emph{inwards}, towards the centre
  (\cref{ch2:fig:flatbend}).
  The tyres grip, so this is static friction. Vertically $F_\mathrm{N} = mg$. Towards the
  centre, $F_\mathrm{f} = \dfrac{mv^2}{r}$. The bend \emph{requires} $\dfrac{mv^2}{r}$, but the
  road can give at most $\mu mg$. If more is required, the car skids towards the outside.
\end{sidebyside}

\begin{formulacard}[notonsheet]{Fastest speed on a flat bend}
  \cardformula{v_\mathrm{max} = \sqrt{\mu g r}}
  \begin{symbols}
    \sym{$v_\mathrm{max}$}{fastest speed without skidding}{\si{m.s^{-1}}}
    \sym{$\mu$}{coefficient of friction (tyres on road)}{no unit}
    \sym{$g$}{gravitational acceleration, \qty{9.8}{m.s^{-2}}}{\si{m.s^{-2}}}
    \sym{$r$}{radius of the bend}{\si{m}}
  \end{symbols}
  \inwords{grippier tyres or a wider bend allow a higher speed. The mass makes no difference.}
  \whentouse{a flat bend where friction is the only sideways force.}
  \whennottouse{on a banked road, or when the friction force is given (use
    $F_\mathrm{f} = \dfrac{mv^2}{r}$).}
  \rearranged{$\mu_\mathrm{min} = \dfrac{v^2}{rg}$ \qquad $r_\mathrm{min} = \dfrac{v^2}{\mu g}$}
  \tensecond{$\mu = 0.50$, $r = \qty{20}{m}$:
    $v_\mathrm{max} = \sqrt{0.50 \times 9.8 \times 20} = \qty{9.9}{m.s^{-1}}$ (\qty{36}{km.h^{-1}}).}
  \derive{friction is the centripetal force, and it has a maximum.
    \begin{align*}
      \frac{mv^2}{r} &\le \mu F_\mathrm{N} = \mu mg && \text{friction needed $\le$ largest friction}\\
      \frac{mv_\mathrm{max}^2}{r} &= \mu mg && \text{at the fastest speed, friction is at its limit}\\
      \frac{v_\mathrm{max}^2}{r} &= \mu g && \text{divide both sides by $m$}\\
      v_\mathrm{max}^2 &= \mu g r && \text{multiply both sides by $r$}\\
      v_\mathrm{max} &= \sqrt{\mu g r} && \text{take the square root}
    \end{align*}}
  \trap{friction is not always $\mu F_\mathrm{N}$: that is the \emph{largest} value. The
    actual friction is whatever is needed, up to that limit.}
\end{formulacard}

\begin{workedexample}{Can the car make the bend?}\label{ch2:we:car}
A \qty{1400}{kg} car drives onto a flat bend of radius \qty{40}{m} at \qty{60}{km.h^{-1}}.
For its tyres, $\mu = 0.80$ on a dry road and $\mu = 0.35$ on a wet road. (a) Find the
friction needed. (b) Can the car take the bend on the dry road? On the wet road?
(c) Find the fastest safe speed on the wet road, in \si{km.h^{-1}}.
\given{$m = \qty{1400}{kg}$, \ $r = \qty{40}{m}$, \ $v = 60 \div 3.6 = \qty{16.667}{m.s^{-1}}$,
  \ $\mu_\mathrm{dry} = 0.80$, \ $\mu_\mathrm{wet} = 0.35$}
\find{the friction needed; the largest friction on each road; $v_\mathrm{max}$ (wet)}
\formula{needed: $F_\mathrm{f} = \dfrac{mv^2}{r}$; \ largest: $\mu mg$; \
  $v_\mathrm{max} = \sqrt{\mu g r}$}
\substitute{\eqlines{\text{needed} &= \frac{\qty{1400}{kg} \times (\qty{16.667}{m.s^{-1}})^2}{\qty{40}{m}}
    = \qty{9.72e3}{N}\\
  \text{largest (dry)} &= 0.80 \times \qty{1400}{kg} \times \qty{9.8}{m.s^{-2}} = \qty{1.10e4}{N}\\
  \text{largest (wet)} &= 0.35 \times \qty{1400}{kg} \times \qty{9.8}{m.s^{-2}} = \qty{4.80e3}{N}\\
  v_\mathrm{max} &= \sqrt{0.35 \times \qty{9.8}{m.s^{-2}} \times \qty{40}{m}} = \qty{11.71}{m.s^{-1}}\\
  &= 11.71 \times 3.6 = \qty{42.2}{km.h^{-1}}}}
\answer{(a) \final{$\qty{9.7e3}{N}$} \ (b) dry: yes ($\qty{1.10e4}{N} > \qty{9.7e3}{N}$);
  wet: no ($\qty{4.8e3}{N} < \qty{9.7e3}{N}$), it skids outwards \ (c) \final{$\qty{42}{km.h^{-1}}$}}
\checkstep{the wet-road limit is below \qty{60}{km.h^{-1}}, which agrees with (b) \cmark.}
\end{workedexample}

% layout QA: split 0.31 (was 0.36): the text column was much taller than the figure
\begin{sidebyside}[0.31]
  \centering
  \begin{tikzpicture}[x=1cm, y=1cm]
    \draw[hscink, line width=1.2pt] (-1.4,-0.2) -- (-1.4,2.6);
    \draw[hscink, line width=1.2pt] (1.4,-0.2) -- (1.4,2.6);
    \draw[hscmuted, line width=0.7pt] (0,2.6) ellipse (1.4 and 0.25);
    \draw[guide] (0,-0.7) -- (0,3.15);
    \node[diagramlabel, anchor=south] at (0,3.15) {axis};
    \draw[->, hscmuted, line width=0.7pt] (-0.5,3.0) arc (200:340:0.55 and 0.18);
    \fill[ground] (-1.3,-0.75) rectangle (1.3,-0.6);
    \draw[hscink, line width=0.7pt] (-1.3,-0.6) -- (1.3,-0.6);
    \node[diagramlabel, anchor=north, text=hscmuted] at (0,-0.8) {floor has dropped};
    \fill[hscbrand!45] (1.15,0.4) rectangle (1.36,1.6);
    \fill[hscbrand!45] (1.255,1.77) circle (0.15);
    \node[diagramlabel, anchor=west] at (1.45,1.75) {rider};
    \node[diagramlabel, anchor=west] at (1.45,2.4) {wall};
    \draw[displacement, <->] (0,0.6) -- (-1.4,0.6) node[pos=0.5, above, diagramlabel] {$r$};
    \coordinate (R) at (1.255,1.05);
    \draw[force] (R) -- (0.1,1.05) node[pos=0.55, above, text=hscforce] {$F_\mathrm{N}$};
    \draw[force] (R) -- (1.255,2.25) node[pos=0.95, left, text=hscforce] {$F_\mathrm{f}$};
    \draw[force] (R) -- (1.255,-0.15) node[pos=0.9, left, text=hscforce] {$mg$};
    \fill[hscink] (R) circle (1.8pt);
  \end{tikzpicture}
  \captionof{figure}{A rotor ride, cut open: the forces on a rider.}\label{ch2:fig:rotor}
\nextside
  In a \term{rotor} ride, a drum spins and the floor drops away, but the riders stay stuck
  to the wall (\cref{ch2:fig:rotor}).
  \begin{itemize}
    \item The wall's normal force is horizontal, towards the axis. It is the centripetal
          force: $F_\mathrm{N} = \dfrac{mv^2}{r}$.
    \item Static friction acts up the wall and balances the weight: $F_\mathrm{f} = mg$.
    \item Riders stay up only if the friction needed ($mg$) is not more than the largest
          friction ($\mu F_\mathrm{N}$). This sets a slowest safe speed (see the card
          below). Mass does not matter.
    \item Faster than this, $F_\mathrm{N}$ and the friction \emph{limit} grow, but the actual
          friction stays equal to $mg$.
  \end{itemize}
\end{sidebyside}

\begin{formulacard}[notonsheet]{Slowest safe speed in a rotor}
  \cardformula{v_\mathrm{min} = \sqrt{\frac{gr}{\mu}}}
  \begin{symbols}
    \sym{$v_\mathrm{min}$}{slowest speed at which riders do not slide down}{\si{m.s^{-1}}}
    \sym{$g$}{gravitational acceleration, \qty{9.8}{m.s^{-2}}}{\si{m.s^{-2}}}
    \sym{$r$}{radius of the drum}{\si{m}}
    \sym{$\mu$}{coefficient of friction (rider on the wall)}{no unit}
  \end{symbols}
  \inwords{a bigger drum or a less grippy wall needs a faster spin. The mass makes no
    difference.}
  \whentouse{a spinning drum (rotor) with the floor taken away, riders pressed on the wall.}
  \whennottouse{a flat bend: there friction is the inward force (use
    $v_\mathrm{max} = \sqrt{\mu g r}$).}
  \rearranged{$\mu_\mathrm{min} = \dfrac{gr}{v^2}$ \qquad longest period
    $T_\mathrm{max} = \dfrac{2\pi r}{v_\mathrm{min}}$}
  \tensecond{$r = \qty{2.0}{m}$, $\mu = 0.50$:
    $v_\mathrm{min} = \sqrt{\dfrac{9.8 \times 2.0}{0.50}} = \qty{6.3}{m.s^{-1}}$.}
  \derive{the wall's push is the centripetal force; friction holds the rider up.
    \begin{align*}
      F_\mathrm{N} &= \frac{mv^2}{r} && \text{horizontal: the wall's push is the net force to the axis}\\
      F_\mathrm{f} &= mg && \text{vertical: the rider does not slide, so the forces balance}\\
      mg &\le \mu F_\mathrm{N} = \mu\frac{mv^2}{r} && \text{friction needed $\le$ largest friction}\\
      mg &= \mu\frac{mv_\mathrm{min}^2}{r} && \text{at the slowest speed, friction is at its limit}\\
      g &= \mu\frac{v_\mathrm{min}^2}{r} && \text{divide both sides by $m$}\\
      gr &= \mu v_\mathrm{min}^2 && \text{multiply both sides by $r$}\\
      v_\mathrm{min}^2 &= \frac{gr}{\mu} && \text{divide both sides by $\mu$}\\
      v_\mathrm{min} &= \sqrt{\frac{gr}{\mu}} && \text{take the square root}
    \end{align*}}
  \trap{in a rotor $F_\mathrm{N} = \dfrac{mv^2}{r}$, not $mg$. Here $\mu$ is \emph{under} the
    line; in $v_\mathrm{max} = \sqrt{\mu g r}$ it is on top.}
\end{formulacard}

\begin{workedexample}{A rotor ride}\label{ch2:we:rotor}
A rotor drum has a radius of \qty{2.5}{m}. The coefficient of friction between a rider and
the wall is 0.50. The rider's mass is \qty{60}{kg}. (a) Find the slowest speed at which the
rider does not slide down. (b) Find the slowest rotation rate, in revolutions per minute
(rpm). (c) At this speed, find the normal force and the friction force on the rider.
\given{$r = \qty{2.5}{m}$, \ $\mu = 0.50$, \ $m = \qty{60}{kg}$}
\find{(a) $v_\mathrm{min}$ \quad (b) turns per minute \quad (c) $F_\mathrm{N}$ and $F_\mathrm{f}$}
\formula{$v_\mathrm{min} = \sqrt{\dfrac{gr}{\mu}}$, \quad $T = \dfrac{2\pi r}{v}$, \quad
  $F_\mathrm{N} = \dfrac{mv^2}{r}$, \quad $F_\mathrm{f} = mg$}
\substitute{\eqlines{v_\mathrm{min} &= \sqrt{\frac{\qty{9.8}{m.s^{-2}} \times \qty{2.5}{m}}{0.50}}
    = \sqrt{\qty{49.0}{m^2.s^{-2}}} = \qty{7.00}{m.s^{-1}}\\
  T_\mathrm{max} &= \frac{2\pi \times \qty{2.5}{m}}{\qty{7.00}{m.s^{-1}}} = \qty{2.244}{s}
    \quad\text{(the longest period)}\\
  \text{rpm} &= \frac{\qty{60}{s}}{\qty{2.244}{s}} = 26.7 \quad\text{(turns in one minute)}\\
  F_\mathrm{N} &= \frac{\qty{60}{kg} \times (\qty{7.00}{m.s^{-1}})^2}{\qty{2.5}{m}} = \qty{1176}{N}\\
  F_\mathrm{f} &= mg = \qty{60}{kg} \times \qty{9.8}{m.s^{-2}} = \qty{588}{N}}}
\answer{(a) \final{$v_\mathrm{min} = \qty{7.0}{m.s^{-1}}$} \ (b) \final{about 27 rpm}
  \ (c) \final{$F_\mathrm{N} = \qty{1.2e3}{N}$} to the axis;
  \final{$F_\mathrm{f} = \qty{5.9e2}{N}$} up the wall}
\checkstep{$\mu F_\mathrm{N} = 0.50 \times \qty{1176}{N} = \qty{588}{N} = mg$: the friction is
  exactly at its limit \cmark. The mass is not needed in (a) or (b) \cmark.}
\end{workedexample}


\begin{mistakes}
\begin{itemize}
  \item \xmark\ Using the string length $L$ as the radius. \cmark\ $r = L\sin\theta$ (centre
        below the pivot).
  \item \xmark\ $F_\mathrm{T} = mg$, or $F_\mathrm{N} = mg$ on a bank. \cmark\ Only the vertical
        \emph{component} equals $mg$.
  \item \xmark\ Drawing $F_\mathrm{N}$ straight up on a banked road. \cmark\ It is at \ang{90}
        to the road. Mark $\theta$ on the diagram before using $\sin$ or $\cos$.
  \item \xmark\ Using $\mu F_\mathrm{N}$ as the actual friction. \cmark\ It is only the largest
        possible friction. In a rotor, $F_\mathrm{f} = mg$.
  \item \xmark\ ``Friction on a turning car points backwards or outwards.'' \cmark\ The
        sideways friction points to the centre of the bend.
\end{itemize}
\end{mistakes}

\begin{examtip}
``Show that $\tan\theta = \dfrac{v^2}{rg}$ for a banked track'' (3 marks, one per step):
\begin{enumerate}
  \item A free-body diagram with only $F_\mathrm{N}$ (at \ang{90} to the road) and $mg$.
  \item The two equations: $F_\mathrm{N}\sin\theta = \dfrac{mv^2}{r}$ and
        $F_\mathrm{N}\cos\theta = mg$.
  \item Divide them to get $\tan\theta = \dfrac{v^2}{rg}$.
\end{enumerate}
``Explain why a car may skid on a wet bend'' (3 marks):
\begin{enumerate}
  \item Friction from the road provides the centripetal force.
  \item A wet road has a smaller $\mu$, so the largest friction, $\mu mg$, is smaller.
  \item If $\dfrac{mv^2}{r} > \mu mg$, the car cannot follow the curve and slides outwards.
\end{enumerate}
\end{examtip}

\begin{quickcheck}
  \item A \qty{0.50}{kg} mass on a string moves in a horizontal circle of radius
        \qty{0.40}{m} at \qty{1.4}{m.s^{-1}} (a conical pendulum). Find the angle of the
        string to the vertical and the tension.
  \item A bend of a cycling track has radius \qty{28}{m} and is banked at \ang{43}. Find
        its design speed in \si{km.h^{-1}}.
  \item A \qty{30}{kg} child on a swing passes the lowest point at \qty{4.0}{m.s^{-1}}. The
        chains are \qty{2.5}{m} long. Find the total tension in the chains there.
  \item A flat bend has radius \qty{65}{m}. Find the fastest safe speed in \si{km.h^{-1}}
        when $\mu = 0.75$ (dry) and when $\mu = 0.10$ (ice). Does the car's mass matter?
  \item A rotor of radius \qty{2.0}{m} turns once every \qty{1.8}{s}. (a) Find the smallest
        coefficient of friction that stops riders sliding down. (b) The rotor then slowly
        spins down. For a rider still stuck to the wall, what happens to the normal force
        and to the friction force?
\end{quickcheck}

%% ==================================================================
% consistency edit: keep the heading with the toolkit below (conditional, no forced break)
\Needspace{14\baselineskip}
\section{Key experiments}\label{ch2:sec:exp}
\covers{M5.IQ2.01}

\begin{toolkit}{Straight-line graphs and the gradient}
The \term{origin} is the point (0, 0) of a graph. If $y = kx$, a graph of $y$ against $x$ is a
straight line through the origin, and its gradient (slope) is $k$. A \term{line of best fit}
is the straight line that passes as close as possible to all the points, with about as many
points above it as below. Write $F_\mathrm{c} = \left(\dfrac{m}{r}\right) v^2$: a graph of
$F_\mathrm{c}$ against $v^2$ should be a straight line with gradient $\dfrac{m}{r}$. Find the
gradient from two points far apart \emph{on the line of best fit} (not two data points):
gradient $= \dfrac{\text{rise}}{\text{run}}$.
\end{toolkit}

\begin{experiment}{Investigating $F_\mathrm{c} = mv^2/r$ with a whirling stopper}
\aim{To investigate how the centripetal force on an object in uniform circular motion
  depends on its speed (Part A), its radius (Part B) and its mass (Part C).}
\apparatus{rubber stopper, strong nylon string, smooth tube (a pen barrel with taped
  ends), slotted masses and hanger, paper clip, metre rule, stopwatch, balance, safety
  glasses.
  \begin{figure}
    \centering
    \begin{tikzpicture}[x=1cm, y=1cm]
    \draw[hscmuted, dashed, line width=0.7pt] (0,0.5) ellipse (3.2 and 0.34);
    \draw[object] (-0.12,-0.75) rectangle (0.12,0.66);
    \draw[hscink, line width=0.8pt] (0,0.66) -- (3.2,0.5);
    \draw[displacement, <->] (0.05,0.98) -- (3.2,0.82) node[pos=0.5, above, diagramlabel] {$L$};
    \fill[hscink] (3.2,0.5) ellipse (0.16 and 0.13);
    \node[diagramlabel, anchor=west] at (3.42,0.5) {stopper, mass $m$};
    \draw[force] (3.05,0.51) -- (2.1,0.56) node[pos=0.5, below, inner sep=1pt, fill=hscexp!4!white, text=hscforce] {$F_\mathrm{T}$};
    \draw[hscink, line width=0.8pt] (0,-0.75) -- (0,-1.35);
    \draw[hscink, fill=hscmuted!40, line width=0.6pt] (-0.12,-0.98) rectangle (0.12,-0.88);
    \node[diagramlabel, anchor=west] at (0.25,-0.93) {paper clip (keeps $L$ fixed)};
    \node[diagramlabel, anchor=west] at (0.25,-0.42) {tube, held in the hand};
    \draw[object] (-0.33,-1.8) rectangle (0.33,-1.35);
    \node[diagramlabel] at (0,-1.575) {$M$};
    \draw[force] (0,-1.8) -- (0,-2.3) node[right, text=hscforce] {$Mg$};
    \node[diagramlabel, anchor=west] at (0.45,-1.575) {hanging mass};
  \end{tikzpicture}
    \caption{Apparatus for the whirling-stopper investigation (side view).}\label{ch2:fig:apparatus}
  \end{figure}}
\method{\textbf{Part A: force and speed.}
\begin{enumerate}
  \item Measure the stopper's mass $m$. Set up as in \cref{ch2:fig:apparatus}, with
        $L = \qty{0.75}{m}$ from the top of the tube to the centre of the stopper.
  \item Whirl the stopper in a horizontal circle. Keep the clip just below the tube, so $L$
        stays the same.
  \item Time 20 turns, three times, and average. Period $T = \text{time} \div 20$.
  \item Repeat for hanging masses $M$ from \qtyrange{0.100}{0.300}{kg}, in steps of
        \qty{0.050}{kg}.
\end{enumerate}
\textbf{Variables in Part A.} The \term{independent variable} (the one you change) is the
hanging mass $M$: it sets the force. The \term{dependent variable} (the one you measure) is the
time for 20 turns: it gives $T$, then $v$. The \term{controlled variables} (kept the same) are
the stopper mass $m$, the length $L$, a horizontal circle and the number of turns timed.

The hanging mass stays still, so the tension equals its weight: $F_\mathrm{T} = Mg$. The
tension provides the centripetal force, so we use $F_\mathrm{c} = Mg$. (The small sag of the
string is dealt with under Validity.) The speed is $v = \dfrac{2\pi L}{T}$.

\textbf{Part B: force and radius.} Keep $M$ (so $F_\mathrm{c}$) and $m$ fixed. Set
$L = 0.40$, 0.55, 0.70, 0.85 and \qty{1.00}{m}, and time 20 turns for each. Plot $v^2$
against $r$ (here $r = L$).

\textbf{Part C: force and mass.} Keep $M$ and $L$ fixed. Use 1, 2, 3 and then 4 identical
stoppers, and time 20 turns for each. Plot $v^2$ against $\dfrac{1}{m}$.}
\results{Part A sample results (made-up but realistic), $m = \qty{0.020}{kg}$,
  $L = \qty{0.75}{m}$. Each time is the mean of three trials, which agreed within
  \qty{0.4}{s}.
  \begin{sidebyside}[0.50]
    \centering\small
    \begin{tabular}{@{}cccc@{}}
      \toprule
      $F_\mathrm{c} = Mg$ & time for & $T$ & $v^2$\\
      (\si{N}) & 20 turns (\si{s}) & (\si{s}) & (\si{m^2.s^{-2}})\\
      \midrule
      0.98 & 15.70 & 0.7850 & 36.0\\
      1.47 & 12.63 & 0.6315 & 55.7\\
      1.96 & 11.10 & 0.5550 & 72.1\\
      2.45 & 9.80  & 0.4900 & 92.5\\
      2.94 & 9.07  & 0.4535 & 108.0\\
      \bottomrule
    \end{tabular}
    \captionof{table}{Part A sample results ($M$ from \qtyrange{0.100}{0.300}{kg}).}\label{ch2:tab:results}
  \nextside
    \centering
    \begin{tikzpicture}
      \begin{axis}[hsc, width=\linewidth, height=0.72\linewidth,
          xmin=0, xmax=120, ymin=0, ymax=3.4,
          xlabel={$v^2$ (\si{m^2.s^{-2}})}, ylabel={$F_\mathrm{c}$ (\si{N})},
          xtick={0,40,...,120}, ytick={0,1,...,3}]
        \addplot[hscseriesB, dashed, line width=1pt, domain=0:120, samples=2] {0.02692*x};
        \draw[hscmuted, densely dotted, line width=0.7pt] (axis cs:0,0) -- (axis cs:110,0)
          -- (axis cs:110,2.96);
        \node[anchor=south, font=\footnotesize, text=hscmuted] at (axis cs:66,0.04) {run $= 110$};
        \node[anchor=east, font=\footnotesize, text=hscmuted] at (axis cs:109,1.3) {rise $= 2.96$};
        \addplot[hscseriesA, only marks, mark=*, mark size=2pt]
          coordinates {(36.0,0.98) (55.7,1.47) (72.1,1.96) (92.5,2.45) (108.0,2.94)};
        \node[anchor=north west, font=\footnotesize, text=hscink, align=left] at (axis cs:4,3.35)
          {best-fit line:\\ gradient $= \dfrac{2.96}{110}$};
      \end{axis}
    \end{tikzpicture}
    \captionof{figure}{$F_\mathrm{c}$ against $v^2$ (Part A).}\label{ch2:fig:graph}
  \end{sidebyside}}
\conclusion{Part A: the graph is a straight line through the origin, so
  $F_\mathrm{c} \propto v^2$ when $m$ and $r$ stay the same. Two points on the best-fit line
  are (0, 0) and (110, 2.96). Gradient $= \dfrac{\text{rise}}{\text{run}}
  = \dfrac{2.96 - 0}{110 - 0} = \qty{0.0269}{kg.m^{-1}}$. Unit:
  $\si{N} \div \si{m^2.s^{-2}} = \si{kg.m.s^{-2}} \div \si{m^2.s^{-2}} = \si{kg.m^{-1}}$. This is
  within 1\,\% of $\dfrac{m}{r} = \dfrac{0.020}{0.75} = \qty{0.0267}{kg.m^{-1}}$, so it supports
  $F_\mathrm{c} = \dfrac{mv^2}{r}$.

  Part B (expected): the theory gives $v^2 = \dfrac{F_\mathrm{c}}{m}\, r$, so $v^2$ against $r$
  should be a straight line through the origin, with gradient $\dfrac{F_\mathrm{c}}{m}$. If it
  is, $F_\mathrm{c} = \dfrac{mv^2}{r}$ holds when $r$ changes. So at a fixed speed, a bigger
  radius needs a smaller force: $F_\mathrm{c} \propto \dfrac{1}{r}$.

  Part C (expected): the theory gives $v^2 = F_\mathrm{c}\, r \times \dfrac{1}{m}$, so $v^2$ against
  $\dfrac{1}{m}$ should be a straight line through the origin, with gradient $F_\mathrm{c}\, r$.
  So at a fixed speed, $F_\mathrm{c} \propto m$.

  Other graphs you may be given: $F_\mathrm{c}$ against $\dfrac{1}{T^2}$ (gradient
  $4\pi^2 mr$), or $T^2$ against $r$ at a fixed force (gradient
  $\dfrac{4\pi^2 m}{F_\mathrm{c}}$). A graph of $F_\mathrm{c}$ against $v$ is a curve, not a
  straight line.}
\significance{It confirms $F_\mathrm{c} = \dfrac{mv^2}{r}$ with a force you can measure (a
  hanging weight). Engineers use the same model to set safe speeds on bends, to design
  banked tracks and rides, and to predict satellite orbits (\cref{ch4:chap}).}
\validity{\begin{itemize}[beginpenalty=10000]
  \item \textbf{Validity} (does it test what it claims?):
        \begin{itemize}
          \item Change one variable at a time.
          \item Keep $L$ fixed with the clip, and keep the circle horizontal.
          \item The hanging mass must stay still, or the tension is not $Mg$.
          \item A smooth tube reduces friction on the string.
          \item The string sags by a small angle $\phi$ below the horizontal. Then the true
                radius is $L\cos\phi$, the true speed is $\dfrac{2\pi L\cos\phi}{T}$, and the
                inward pull is the horizontal part of the tension, $F_\mathrm{T}\cos\phi$. Put
                these into $F_\mathrm{c} = \dfrac{mv^2}{r}$:
                \[ F_\mathrm{T}\cos\phi = \frac{m(2\pi L\cos\phi)^2}{T^2 \times L\cos\phi}
                   = \frac{4\pi^2 mL\cos\phi}{T^2}. \]
                Divide both sides by $\cos\phi$: $F_\mathrm{T} = \dfrac{4\pi^2 mL}{T^2}$. Using
                $L$ for $r$ and $v = \dfrac{2\pi L}{T}$ gives the same result, so the sag
                cancels out.
        \end{itemize}
  \item \textbf{Reliability} (do repeats agree?): three trials per mass, averaged; five or
        more masses; a line of best fit.
  \item \textbf{Accuracy} (how close to the true value?): reaction time is about
        \qty{0.2}{s}. That is 44\,\% of one short period (\qty{0.45}{s}), but only 2\,\% of
        20 periods (about \qty{9}{s}). So time 20 turns, or film with a phone. Measure $L$ to
        the centre of the stopper.
  \item \textbf{In the exam,} ``human error'' earns nothing: name the error, say how it
        affects the result, and give a specific fix. Safety: glasses, tight knots, clear space.
\end{itemize}}
\end{experiment}

% Added by the syllabus coverage audit (M5.IQ2.01 is a required investigation, but this section
% had no exam tip and no Quick check, and Parts B and C had no data to analyse).
\begin{examtip}
Three kinds of question are common for this investigation.
\begin{itemize}
  \item ``Outline a procedure'' (3--4 marks). Say what you change and how: the hanging mass
        $M$ sets $F_\mathrm{c} = Mg$. Say what you measure and how: time 20 turns, then find
        $T$ and $v = 2\pi r/T$. Say what you keep the same: $m$, and $r$ (with the
        paper clip). Repeat each trial, and plot a graph.
  \item ``Explain how the results support $F_\mathrm{c} = mv^2/r$'' (2--3 marks).
        Choose the graph that should be straight: $F_\mathrm{c}$ against $v^2$, $v^2$ against
        $r$, or $v^2$ against $1/m$. Say that it is a straight line through the
        origin. Then compare its gradient with the predicted value, as a percentage
        difference.
  \item ``Evaluate the method'' (4--5 marks). For each problem, name it, give its effect on
        the result, and give a specific fix (see ``Validity, reliability and accuracy'' in the
        experiment above, and \cref{pq:stopper}).
\end{itemize}
\end{examtip}

\begin{quickcheck}
  \item Part B: a student keeps $M = \qty{0.150}{kg}$ and uses a \qty{0.025}{kg} stopper. Her
        graph of $v^2$ against $r$ is a straight line through the origin, with gradient
        \qty{57.5}{m.s^{-2}}. (a) What does a straight line through the origin show? (b) Use
        the gradient to find $F_\mathrm{c}$. Compare it with $Mg$ as a percentage difference.
        (c) Use the line to predict the time for 20 turns when $r = \qty{0.80}{m}$.
  \item Part C: a student keeps $M$ and $L$ fixed. She uses 1, 2, 3 and then 4 identical
        stoppers. (a) Name the independent variable, the dependent variable and two
        controlled variables. (b) Which graph should be a straight line through the origin?
        What should its gradient be? (c) With one stopper the period is \qty{0.60}{s}.
        Predict the period with two stoppers.
  \item Use \cref{ch2:tab:results}. When $F_\mathrm{c}$ doubles from \qty{0.98}{N} to
        \qty{1.96}{N}, by what factor does $v$ change? Explain why a graph of $F_\mathrm{c}$
        against $v$ is a curve, but a graph of $F_\mathrm{c}$ against $v^2$ is a straight
        line.
\end{quickcheck}

%% ==================================================================
% Chapter summary (added by the whole-book consistency edit: "How to use this guide" says
% that every chapter ends with a short summary). Every statement is taught above.
\begin{summarybox}[unbreakable]{\Cref{ch2:chap} in five lines}
\begin{itemize}
  \item An object moving in a circle at a constant speed is accelerating, because its
        direction changes. The acceleration $a_\mathrm{c} = \frac{v^2}{r}$ points to the centre.
  \item The net force towards the centre is $F_\mathrm{c} = \frac{mv^2}{r}$. It is not a new force:
        always name the real force (tension, friction, a normal force or gravity).
  \item Method: draw only real forces. Split them into a part towards the centre (total
        $\frac{mv^2}{r}$) and a vertical part (total zero for a horizontal circle).
  \item Conical pendulum, or banked road at its design speed: $\tan\theta = \frac{v^2}{rg}$. Flat bend:
        $v_\mathrm{max} = \sqrt{\mu g r}$. Rotor: $v_\mathrm{min} = \sqrt{\frac{gr}{\mu}}$.
  \item Whirling stopper: $F_\mathrm{c} = Mg$, and a graph of $F_\mathrm{c}$ against $v^2$ is a
        straight line through the origin, with gradient $\frac{m}{r}$.
\end{itemize}
\end{summarybox}
